\section{Bài mở đường nói gì, và từ chối nói gì}
\label{sec:ch17-bai-mo-duong}

\index{chuỗi suy luận!bài mở đường tự đặt giới hạn|(}%
Mục trước khép lại ở chỗ đối tượng đã đổi một lần kể từ bài mở đường. Mục này
quay về chính bài ấy, vì trước khi hỏi ngành đã đọc nó thế nào thì phải biết nó
viết gì. Bài này chương~\ref{ch:mo-hinh-can-giai-thich} đã đọc một lần, cho câu
hỏi chuỗi suy luận là gì; ở đây nó được đọc cho một câu hỏi khác, là bài có
tuyên bố chuỗi suy luận giải thích được mô hình hay không.

\index{chuỗi suy luận!bốn tính chất hấp dẫn}%
Bài nêu bốn tính chất hấp dẫn của cách nhắc này, ở trang 3, và tính chất thứ hai
là tính chất duy nhất nói về việc hiểu mô hình~\autocite{p06cot}. Nó được viết ra
như sau, kể cả phần trong ngoặc:

\begin{quote}
Thứ hai, một chuỗi suy luận cung cấp một khung nhìn diễn giải được vào hành vi
của mô hình, gợi ra \emph{cách mà nó có thể đã} đi tới một đáp án cụ thể và cho
cơ hội gỡ lỗi chỗ đường suy luận đi sai (\emph{mặc dù việc mô tả đầy đủ những
tính toán của một mô hình nâng đỡ cho một đáp án vẫn là một câu hỏi mở}).%
\footnote{Nhấn mạnh là của sách. Nguyên văn: \enquote{Second, a chain of thought
provides an interpretable window into the behavior of the model, suggesting how
it might have arrived at a particular answer and providing opportunities to debug
where the reasoning path went wrong (although fully characterizing a model's
computations that support an answer remains an open question).}}
\end{quote}

Hai chỗ rào trong một câu, và cả hai đều mang tải. \enquote{Có thể đã} nói rằng
chuỗi gợi ra một lời kể chứ không xác lập nó. Phần trong ngoặc thì nói thẳng
rằng việc đối chiếu chuỗi với những tính toán thật là chuyện chưa ai làm được,
tức là đúng vế phải của định nghĩa~\ref{def:ch01-do-trung-thuc}. Sách trích cả
câu chứ không trích mệnh đề đầu, vì mệnh đề đầu đứng một mình chính là cách đọc
mà mục~\ref{sec:ch17-nganh-doc-nguoc-lai} sẽ tìm thấy.

\index{chuỗi suy luận!bài từ chối gọi nó là lời giải thích}%
Bài còn từ chối một chuyện nữa, và lần này là từ chối một cái tên. Ở trang 3, khi
giới thiệu ví dụ đầu tiên, bài viết rằng chuỗi suy luận trong trường hợp ấy trông
giống một lời giải và có thể được đọc như một lời giải, nhưng các tác giả vẫn
chọn gọi nó là chuỗi suy luận để giữ được ý rằng nó \emph{bắt chước} một quá
trình nghĩ theo từng bước~\autocite{p06cot}. Động từ mà bài chọn cho quan hệ giữa
văn bản và quá trình là \enquote{bắt chước}, chứ không phải một động từ nào nói
rằng văn bản ghi lại quá trình.

\index{chuỗi suy luận!khả năng diễn giải là hệ quả phụ}%
Câu dứt khoát nhất thì nằm ở mục C.2 của phụ lục, trang 24, chỗ bài tách mình
khỏi một dòng nghiên cứu khác. Dòng ấy là các lời giải thích bằng ngôn ngữ tự
nhiên, với mục đích cải thiện
\tn{khả năng diễn giải}{interpretability} của mạng nơ-ron. Bài phân biệt mình với
dòng ấy bằng hai điểm~\autocite{p06cot}. Điểm thứ nhất là thứ tự: lời giải thích
bằng ngôn ngữ tự nhiên được sinh ra \emph{hoặc cùng lúc với, hoặc sau} tiên đoán
cuối cùng, còn chuỗi suy luận thì đi trước đáp án. Vế \enquote{cùng lúc với} phải
giữ nguyên, vì mục~\ref{sec:ch17-loi-giai-thich-tu-viet} vừa dựng hệ quả thứ hai
của nó trên chính chỗ đồng thời ấy. Điểm thứ hai là mục đích, và bài phát biểu nó
như sau:

\begin{quote}
Và trong khi lời giải thích bằng ngôn ngữ tự nhiên chủ yếu nhắm tới việc cải
thiện khả năng diễn giải của mạng nơ-ron, mục tiêu của cách nhắc theo chuỗi suy
luận là cho phép mô hình phân rã các nhiệm vụ suy luận nhiều chặng thành nhiều
bước; \emph{khả năng diễn giải chỉ là một hệ quả phụ}.%
\footnote{Nhấn mạnh là của sách. Nguyên văn của vế đầu: \enquote{And while NLE
aims mostly to improve neural network interpretability (Rajagopal et al., 2021),
the goal of chain-of-thought prompting is to allow models to decompose multi-hop
reasoning tasks into multiple steps}; bài nối thẳng sang vế sau bằng một gạch
ngang dài không có khoảng trắng hai bên, rồi viết
\enquote{interpretability is just a side effect.}}
\end{quote}

Chỗ đứng của câu này đáng ghi lại, vì nó quyết định câu ấy nặng đến đâu. Nó không
phải một lời khiêm tốn đặt trong phần hạn chế, mà là câu bài dùng để nói mình
khác dòng kia ở chỗ nào. Nghĩa là bài không chỉ tránh tuyên bố về khả năng diễn
giải; nó lấy việc không nhắm tới khả năng diễn giải làm dấu hiệu nhận dạng của
chính đóng góp mình.

\index{chuỗi suy luận!hai hạn chế bài tự ghi}%
Phần hạn chế nêu bốn điểm và hai trong số đó thuộc về chương này. Điểm thứ nhất
viết rằng tuy chuỗi suy luận mô phỏng quá trình nghĩ của người suy luận, điều ấy
không trả lời câu hỏi mạng nơ-ron có thật sự đang \enquote{suy luận} hay không,
và các tác giả để ngỏ câu hỏi ấy~\autocite{p06cot}. Điểm thứ ba viết rằng không
có gì bảo đảm các đường suy luận là đúng, nên một chuỗi được sinh ra có thể dẫn
tới cả đáp án đúng lẫn đáp án sai~\autocite{p06cot}; câu này
mục~\ref{sec:ch02-chuoi-suy-luan} đã dẫn.

\index{chuỗi suy luận!một chữ vắng mặt}%
Còn một dữ kiện nữa và nó là dữ kiện mạnh nhất trong mục này. Rà toàn văn 43
trang, sau khi nối các chỗ ngắt dòng có gạch nối và chuẩn hóa khoảng trắng, chữ
\enquote{faithfulness} xuất hiện đúng hai lần và chữ \enquote{unfaithful} không
lần nào. Lần thứ nhất nằm trong danh mục tài liệu, bên trong nhan đề một bài
được trích. Lần thứ hai nằm ở trang 16, trong một dãy đặt trong ngoặc liệt kê
các năng lực trồi lên theo quy mô: hiểu ngữ nghĩa, ánh xạ ký hiệu, bám chủ đề,
khả năng số học, và \enquote{faithfulness}~\autocite{p06cot}. Nghĩa ở đó là một năng lực mà
mô hình được cho là có, không phải một quan hệ giữa lời giải thích và tính toán,
nên nó không phải nghĩa của định nghĩa~\ref{def:ch01-do-trung-thuc}.

\index{chuỗi suy luận!bài mở đường tự đặt giới hạn|)}%
Vậy bài đưa chuỗi suy luận vào ngành đã nói ba lần rằng nó không tuyên bố chuỗi
ấy giải thích mô hình: một lần bằng phần trong ngoặc của tính chất thứ hai, một
lần bằng việc chọn chữ \enquote{bắt chước} thay cho chữ \enquote{lời giải
thích}, và một lần bằng câu gọi khả năng diễn giải là hệ quả phụ. Nó cũng chưa
một lần dùng chữ chỉ tính chất mà sách này bàn. Điều mục sau đọc là chuyện xảy ra
sau đó.
